% Generated from locale/tr/content/normal-modal-logic/frame-definability/properties-accessibility.tex; do not hand-edit.


\olfileid[tr]{nml}{frd}{acc}

\olsection{Erişilebilirlik Bağıntılarının Özellikleri}

Birçok modal !!{formula}{}ün erişilebilirlik bağıntısının basit, hatta
tanıdık özelliklerini nitelendirdiği ortaya çıkar. Bunun bir yönü şudur:
belirli bir özelliğe sahip her model, karşılık gelen bir !!{formula}{}ü (ve
onun bütün yerine koyma örneklerini) doğru kılar. Erişilebilirlik
bağıntılarının beş klasik türüyle ve doğruluğunu güvence altına aldıkları
formüllerle başlayalım.

\begin{thm}\ollabel{thm:soundschemas}
  $\mModel{M}=\tuple{W,R,V}$ bir model olsun. $R$, \olref{tab:five}
  tablosunun solundaki özelliğe sahipse, sağdaki !!{formula}{}ün her örneği
  $\mModel{M}$ içinde doğrudur.
\end{thm}

\begin{table}[t]
    \begin{tabular}{| l || l |}
      \hline
      {\emph{$R$ \dots ise}} & {\emph{\dots $\mModel{M}$ içinde doğrudur:}} \\
      \hline \hline
      \emph{seri}: $\forall u \exists v Ruv$ & \hbox to.43\textwidth
           {$\Box p \lif \Diamond p$ \hfill (\Ax{D})} \\
      \hline
      \emph{yansımalı}: $\forall w Rww$
      & \hbox to.43\textwidth
          {$\Box p \lif p$ \hfill (\Ax{T})} \\
      \hline
      \emph{simetrik}: &  \hbox to .43\textwidth
           {$p \lif \Box\Diamond p$ \hfill (\Ax{B})} \\
           $\forall u\forall v(Ruv \lif Rvu)$ & \\
      \hline
      \emph{geçişli}: & \hbox to.43\textwidth
           {$\Box p \lif \Box \Box p$ \hfill (\Ax{4})} \\
      $\forall u \forall v \forall w ((Ruv \land Rvw) \lif Ruw)$ & \\
      \hline
      \emph{Öklidyen}: & \hbox to .43\textwidth
        {$\Diamond p \lif \Box\Diamond p$ \hfill (\Ax{5})} \\
      $\forall w \forall u \forall v ((Rwu \land Rwv) \lif Ruv)$ &\\
      \hline
    \end{tabular}
    \caption{Beş karşılıklılık olgusu.}
    \ollabel{tab:five}
  \end{table}


\begin{proof}
  \Ax{B} durumunu ele alalım: şemanın bir modelde doğru olduğunu göstermek
  için bütün örneklerinin modelin bütün dünyalarında doğru olduğunu
  göstermeliyiz. Öyleyse \Ax{B}'nin belirli bir $!A\lif\Box\Diamond !A$
  örneğini ve gelişigüzel bir $w\in W$ dünyasını alalım. $\Box\Diamond !A$'nın
  $w$'de doğru olduğunu göstermek üzere önbileşen $!A$'nın $w$'de doğru
  olduğunu varsayalım. Böylece $\Diamond !A$'nın $w$'den erişilebilen bütün
  $w'$ dünyalarında doğru olduğunu göstermemiz gerekir. $Rww'$ olan herhangi
  bir $w'$ için simetri varsayımından $Rw'w$ de elde edilir (bkz.
  \olref{fig:Bsymm}). $\mSat{M}{!A}[w]$ olduğundan
  $\mSat{M}{\Diamond !A}[w']$ olur. $w'$, $Rww'$ olan gelişigüzel bir dünya
  olduğundan $\mSat{M}{\Box\Diamond!A}[w]$ sonucuna varırız.

  Öteki durumları alıştırma olarak bırakıyoruz.
\end{proof}

\begin{prob}
  \olref[nml][frd][acc]{thm:soundschemas} teoreminin ispatını tamamlayın.
\end{prob}

\begin{figure}
  \begin{center}
    \begin{tikzpicture}[modal]
      \node[world] (w1) [label={below:$\mSat{{}}{\formula{A}}$\\
          $\mSat{{}}{\Box\Diamond \formula{A}}$}] {$w$};
      \node[world] (w2) [label=below:{$\mSat{{}}{\Diamond \formula{A}}$},
        right=of w1] {$w'$};
      \draw[->, bend left] (w1) to (w2);
      \draw[->, bend left] (w2) to (w1);
    \end{tikzpicture}
  \end{center}
\caption{Simetriden çıkarılan akıl yürütme.}
\ollabel{fig:Bsymm}
\end{figure}

\olref{thm:soundschemas} teoremindeki gerektirmelerin terslerinin geçerli
olmadığına dikkat edin: bir modelin bir şemayı doğrulamasından, o modelin
erişilebilirlik bağıntısının karşılık gelen özelliğe sahip olduğu çıkmaz.
\Ax{T} ve yansımalı olmayan modeller durumunda, \Ax{T}'nin kendisinin
başarısız olduğu bir model örneği vermek kolaydır: $W=\{w\}$,
$R=\emptyset$ ve $V(p)=\emptyset$ olsun. Bu durumda $R$ yansımalı değildir, fakat
$\mSat{M}{\Box p}[w]$ ve $\mSat/{M}{p}[w]$ olur. Ancak burada \Ax{T}'nin
$\mModel{M}$ içinde başarısız olan yalnızca tek bir örneği vardır; örneğin
$\Box\lnot p\lif\lnot p$ gibi başka örnekler doğrudur. \Ax{T}'nin
\emph{her yerine koyma örneğinin} $\mModel{M}$ içinde doğru, ama
$\mModel{M}$'nin yansımalı olmadığı örnekler vermek daha zordur. Yine de bu
tür modeller de vardır:

\begin{prop}\ollabel{prop:reflexive}
  $\mModel{M}=\tuple{W,R,V}$, $W=\{u,v\}$ ve
  $R=\{(u,v),(v,u)\}$ olan bir model olsun.
  Bütün $p$'ler için $u\in V(p)\Leftrightarrow v\in V(p)$ olduğunu
  varsayalım. O zaman:
  \begin{enumerate}
  \item Bütün $!A$'lar için $\mSat{M}{!A}[u]$ ancak ve ancak
    $\mSat{M}{!A}[v]$ olur ($!A$ üzerinde tümevarım kullanın).
  \item \Ax{T}'nin her örneği $\mModel{M}$ içinde doğrudur.
  \end{enumerate}
  $\mModel{M}$ yansımalı olmadığından (hatta \emph{yansımasızdır}),
  \olref{thm:soundschemas} teoreminin tersi \Ax{T} durumunda başarısızdır
  (\olref{thm:soundschemas} içinde anılan öteki şemaların bazıları---ama
  hepsi değil---için de benzer akıl yürütmeler verilebilir).
\end{prop}

\begin{prob}
  \olref[nml][frd][acc]{prop:reflexive} önermesindeki iddiaları ispatlayın.
\end{prob}

Beş klasik !!{formula} olan \Ax{D}, \Ax{T}, \Ax{B}, \Ax{4} ve \Ax{5}
üzerinde duracak olsak da erişilebilirlik bağıntılarının birkaç başka
özelliğini \olref{tab:anotherfive} tablosunda kaydediyoruz. Her dünyadan en
fazla bir dünyaya erişilebiliyorsa erişilebilirlik bağıntısı $R$
kısmen fonksiyoneldir. Her dünyadan tam olarak bir dünyaya erişilebiliyorsa
ona fonksiyonel deriz. (Dolayısıyla fonksiyonel bağıntılar tam olarak hem
seri hem kısmen fonksiyonel olanlardır.) Erişilebilirlik bağıntısı bir
(kısmi) fonksiyon gibi işlediği için bunlara ``fonksiyonel'' denir. $Ruv$
olduğunda $u$ ile $v$ ``arasında'' bir $w$ varsa bağıntı zayıf yoğundur.
Bu bakımdan zayıf yoğun bağıntılar bir anlamda geçişli bağıntıların tersidir:
geçişli bir bağıntıda $u$'dan $v$'ye $w$ üzerinden dolanarak
ulaşabiliyorsanız $u$'dan $v$'ye doğrudan da ulaşabilirsiniz; zayıf yoğun
bir bağıntıda ise $u$'dan $v$'ye doğrudan ulaşabiliyorsanız bir $w$
üzerinden dolanarak da ulaşabilirsiniz. Bir $w$ dünyasından $u$ ve $v$
dünyalarına ulaşabildiğiniz her durumda hem $u$'dan hem $v$'den tek bir $t$
dünyasına ulaşabiliyorsanız bağıntı zayıf yönlendirilmiştir---buna bazen
``elmas özelliği'' ya da ``birleşme'' denir.

\begin{table}[t]
    \begin{tabular}{| p{.50\textwidth} || p{.43\textwidth} |}
      \hline
      {\emph{$R$ \dots ise}} & {\emph{\dots $\mModel{M}$ içinde doğrudur:}} \\
      \hline \hline
      \emph{kısmen fonksiyonel}: &
      \multirow{2}{*}{$\Diamond p \lif \Box p$} \\
      $ \forall w \forall u \forall v ((Rwu \land Rwv) \lif u =v )$ & \\
      \hline
      \multirow{2}{*}{\emph{fonksiyonel}: $\forall w \exists v \forall u(Rwu \liff \eq[u][v]) $}
      & \multirow{2}{*}{$\Diamond p \liff \Box p$} \\
      & \\
      \hline
      \emph{zayıf yoğun}: &  \multirow{2}{*}{$\Box \Box p \lif
        \Box p$} \\
      $\forall u\forall v(Ruv \lif \exists w(Ruw \land Rwv))$ & \\
      \hline
      \emph{zayıf bağlantılı}: &
      \multirow{3}{*}{\hbox to.43\textwidth{$
        \begin{array}{@{}l@{}}
          \Box ((p \land \Box p) \lif q) \lor {} \\
          \qquad \Box ((q \land \Box q) \lif p)
        \end{array}$ \hfill (\Ax{L})}
      } \\
      $\forall w \forall u \forall v ((Rwu \land Rwv) \lif {}$ &\\
      \qquad $(Ruv \lor u=v \lor Rvu))$ & \\
      \hline
      \emph{zayıf yönlendirilmiş}: & \multirow{3}{*}{\hbox to .43\textwidth
        {$\Diamond\Box p \lif \Box\Diamond p$\hfill (\Ax{G})}} \\
      $\forall w \forall u \forall v ((Rwu \land Rwv) \lif {}$ & \\
      \qquad $\exists t (Rut \land Rvt))$ & \\
      \hline
    \end{tabular}
    \caption{Beş başka karşılıklılık olgusu.}
    \ollabel{tab:anotherfive}
  \end{table}

\begin{prob}
  $\mModel{M}=\tuple{W,R,V}$ bir model olsun. $R$,
  \olref[nml][frd][acc]{tab:anotherfive} tablosunun solundaki özellikleri
  sağlıyorsa, sağda karşılık gelen formülün her örneğinin $\mModel{M}$
  içinde doğru olduğunu gösterin.
\end{prob}



